\documentclass[11pt]{article}
\usepackage{amsmath, amssymb, amsthm}
\usepackage{graphicx}
\usepackage{booktabs}
\usepackage{hyperref}
\usepackage{algorithm}
\usepackage{algpseudocode}
\usepackage{geometry}
\geometry{margin=1in}

\title{The Simplex Method in Modern Optimization: \\A Comparative Study with Gradient Descent}
\author{Ryan Saint-Yves}
\date{Final Project Report \\ 2025}

\begin{document}
\maketitle

\section{Introduction}

Linear programming (LP) is a foundational problem class in optimization, appearing in resource allocation, operations research, finance, and machine learning. This project explores two contrasting algorithmic approaches to optimization on convex polytopes: the Simplex Method and Gradient Descent. The Simplex Method navigates the vertices of the feasible polytope, while gradient descent (adapted with barrier functions) follows continuous descent directions through the interior.


\subsection{Project Motivation}
The Simplex method is a wonderful optimization algorithm for linear programs. Both the Simplex Method and Gradient Descent have reputations as practical algorithms for optimization. Despite its worst-case exponential complexity, Simplex performs efficiently on most practical LPs. Gradient Descent, on the other hand, is simple, scalable, and forms the basis of modern large-scale optimization, though it requires modifications such as penalty or barrier terms to handle polyhedral constraints.

This report aims to answer fundamental questions about the Simplex Algorithm by comparing it to Gradient Descent:
\begin{itemize}
    \item \textbf{What does it do?} How does it approach the optimization problem geometrically and algorithmically?
    \item \textbf{How fast does it run?} What is its computational performance across different problem scales?
    \item \textbf{What are its strengths and weaknesses?} Under what conditions does each algorithm excel or fail?
    \item \textbf{Why is it useful?} What practical insights guide algorithm selection?
\end{itemize}

\subsection{Project Scope}
This report presents:
\begin{itemize}
    \item A theoretical comparison of Simplex vs. Gradient Descent, including feasibility, convergence, and geometric intuition.
    \item A systematic empirical study across small, medium, and large synthetic LPs with varying structural properties.
    \item Analysis of degeneracy, cycling, sensitivity to scaling, and numerical stability issues.
    \item Practical recommendations for algorithm selection based on problem characteristics.
    \item Other algorithms derived from the Simplex method
\end{itemize}


%%%%%%%%%%%%%%%%%%%%%%%%%%%%%%%%%%%%%%%%%%%%%%%%%%%%%%%%%%%%%%%%%%%%%%%%%%%%%%%%%%%%%%%%%%%%%%%%
%%%%%%%%%%%%%%%%%%%%%%%%%%%%%%%%%%%%%%%%%%%%%%%%%%%%%%%%%%%%%%%%%%%%%%%%%%%%%%%%%%%%%%%%%%%%%%%%
%SECTION 2
%%%%%%%%%%%%%%%%%%%%%%%%%%%%%%%%%%%%%%%%%%%%%%%%%%%%%%%%%%%%%%%%%%%%%%%%%%%%%%%%%%%%%%%%%%%%%%%%
%%%%%%%%%%%%%%%%%%%%%%%%%%%%%%%%%%%%%%%%%%%%%%%%%%%%%%%%%%%%%%%%%%%%%%%%%%%%%%%%%%%%%%%%%%%%%%%%


\section{Background}
This section introduces the mathematical formulation of LPs and summarizes the key ideas behind the Simplex Method and Gradient Descent (with penalty or barrier approaches).

\subsection{The Linear Programming Problem}
We consider the standard-form linear program
\begin{align}
    \text{maximize} \quad & c^\top x \\
    \text{subject to} \quad & Ax \le b, \\
    & x \ge 0,
\end{align}
where \(A \in \mathbb{R}^{m \times n}\), \(b \in \mathbb{R}^m\), and \(c \in \mathbb{R}^n\).

To build geometric intuition, first consider a linear program with only two variables.  
A constraint of the form
\[
ax_1 + bx_2 \le c
\]
defines a half-space in \(\mathbb{R}^2\): the region on one side of the line \(ax_1 + bx_2 = c\).  
Adding more constraints produces the intersection of several half-spaces, which is always a convex polygon (in two dimensions) or, more generally, a convex polytope.

\begin{figure}[h]
    \centering
    \includegraphics[width=0.18\linewidth]{images/polytope_equations.jpg}
    \includegraphics[width=0.3\linewidth]{images/polytope.jpg}
    \caption{2D constraints and corresponding polytope}
    \label{fig:1}
\end{figure}

Because the feasible region is convex and we have a clear objective function, it is natural to imagine applying gradient-based methods.  
However, in linear programming, the linearity of the objective gives us a much stronger insight:  
\emph{the optimal solution, if it exists and is bounded, always occurs at a vertex of the feasible polytope}.

To see this, consider a linear objective such as
\[
c_1 x_1 + c_2 x_2 = b.
\]
This objective function is represented by a straight line. Increasing the objective corresponds to shifting this line in the direction of its gradient while keeping it parallel to itself. As we translate the line, it “slides’’ across the feasible region until it can move no farther without leaving the polytope. At that final position, the line touches the feasible region at a single point (or possibly along an entire edge in degenerate cases). In the typical case, this point is a vertex of the polytope.

\begin{figure}[h]
    \centering
    \includegraphics[width=0.4\linewidth]{images/sliding_objective.jpg}
    \caption{Sliding the objective function until it reaches an optimal vertex}
    \label{fig:2}
\end{figure}

This geometric fact motivates the design of the simplex method: rather than searching the interior of the feasible region, we only need to examine its vertices and move along edges from one vertex to another, while improving the objective value.

\subsubsection{Linear Program Example}

One example of this is resource allocation. We may have 10 different resources labeled \(x_1\) through \(x_\text{10}\), maybe wood, metal, cotton etc. And say we want to maximize our value, where each resource is worth different amounts (I.e \$ 10 per unit of wood, \$ 20 per unit of metal...). And finally, we have limitations; we can only physically hold so much of each resource due to various factors. This provides us with constraints like:
\begin{itemize}
    \item \(3x_\text{wood} + 6x_\text{cotton} \le 10\)
    \item \(x_\text{wood} + x_\text{metal} + x_\text{cotton} \le 100\)
\end{itemize}

Now we can set this up using a matrix \(A\) and a vector \(b\). Each row in \(A\) represents the coefficients in our constraint, while the rows of \(b\) are the values that our constraints must satisfy. This way, we can represent as many constraints on as many variables as we want. We then introduce the vector \(c\) to represent how much each variable contributes to the value we want to maximize.


\subsection{The Simplex Method}

We now outline how the simplex method solves a linear program by moving from vertex to vertex of the feasible polytope and improving the objective at each step.

The first step is to convert all inequality constraints into equalities by introducing \emph{slack variables}. For instance, the constraint
\[
ax_1 + bx_2 \le c
\]
becomes
\[
ax_1 + bx_2 + s_1 = c.
\]
Each slack variable represents the unused capacity in its constraint. After adding slacks, we obtain an initial \emph{basic feasible solution} (BFS) by taking all slack variables as basic and the original variables as nonbasic. We will come back to this BFS shortly.

With this BFS, we form the initial simplex tableau. The tableau contains:
\begin{itemize}
    \item one row for each constraint (each basic variable)
    \item one row for the objective function and its reduced costs (how much a variable improves our objective)
    \item one column for each decision variable and slack variable;
    \item a \emph{solution column} giving the current values of the basic variables.
\end{itemize}

A crucial idea is that nonbasic variables are fixed to zero, while each basic variable appears in exactly one row and can be read directly from the solution column. The coefficients in any row can be interpreted as the tradeoff between variables: how much one variable changes when another is increased by one unit.

For example, consider the row
\[
3x + s = 1.
\]
Here \(s\) is basic and \(x\) is nonbasic. Setting \(x = 0\) gives the BFS value \(s = 1\). Increasing \(x\) forces
\[
s = 1 - 3x,
\]
so producing one unit of \(x\) “costs’’ 3 units of slack. This illustrates how simplex tracks opportunity costs and explains why the pivot operation is necessary.

Another note here: To maximize \(Z = ax_1 + bx_2 ...\) we have the initial solution:
\(Z - ax_1 - bx_2 -... = 0\), with \(x_i = 0 \space \forall i\). As you can see, our objective row initially has negative coefficients. Keep this in mind as we move on to discuss pivot selection, the driving mechanism behind the Simplex Method.

\begin{figure}[h]
    \centering
    \includegraphics[width=1\linewidth]{images/initial_tableau.jpg}
    \caption{Initial tableau}
    \label{fig:tableau}
\end{figure}

\subsubsection*{Pivot Selection}
To move to a better vertex, the simplex method chooses:
\begin{enumerate}
    \item an \emph{entering variable}, typically the nonbasic variable with the most negative reduced cost (for maximization). Increasing this variable most improves the objective.
    \item a \emph{leaving variable}, determined using the \emph{minimum ratio test}. For each constraint row \(i\), compute
    \[
    \frac{A_{i,\text{solution}}}{A_{i,\text{entering}}}
    \]
    whenever \(A_{i,\text{entering}} > 0\). This gives the maximum amount the entering variable can increase before violating a constraint. The smallest ratio identifies the tightest constraint, whose basic variable must leave.
\end{enumerate}

\subsubsection*{Pivot Operation}
Once the variables are selected, we make a pivot at the entry \(A_{i,j}\) corresponding to the chosen row and column. This operation reflects in our table the effect of choosing the entering variable on the other constraints and the objective function. 

The pivot operation is identical to Gaussian elimination:
\begin{enumerate}
    \item Divide the pivot row by the pivot element so that the entering variable has coefficient 1.
    \item Use row operations on all other rows (including the objective row) to make all other entries in that column zero.
\end{enumerate}

\begin{figure}[h]
    \centering
    \includegraphics[width=1\linewidth]{images/pivot_tableau.jpg}
    \caption{Choosing pivot in tableau}
    \label{fig:tableau}
\end{figure}

This updates the tableau to reflect a new BFS at a new vertex of the feasible polytope. The entering variable becomes basic, and the leaving variable becomes nonbasic. Since we had to take x amount of our entering variable, we have to reflect that change on the other constraints. First, we normalize for simplicity, then we remove that variable from other constraints. From the example table, \(x_1\) enters and \(s_1\) leaves. But since we've taken \(x_1\), this must be reflected in the constraint corresponding to \(s_4\), so we subtract the \(x_1\) row from the \(s_4\) row to reflect the changes needed to maintain feasibility while adding \(x_1\) to our basis.

\begin{figure}[h]
    \centering
    \includegraphics[width=1\linewidth]{images/updated.jpg}
    \caption{Update tableau}
    \small{Our objective row now reflects our current objective value of 3. A negative coefficient indicates we are not yet at our optimal solution}
    \label{fig:tableau}
\end{figure}


We repeat this process—selecting entering and leaving variables, then pivoting—until no reduced cost indicates further improvement. At that point, the algorithm has reached an optimal vertex.

\subsubsection{Key Properties of the Simplex Method}
The Simplex Method has several important theoretical and practical properties:
\begin{itemize}
    \item It maintains feasibility at every iteration and therefore always stays on vertices of the feasible polytope.
    \item It terminates at an optimal vertex (or detects unboundedness or infeasibility).
    \item It moves along edges of the polytope, changing only one basic variable at a time.
    \item It selects entering and leaving variables using reduced costs and pivot rules.
    \item Degeneracy can cause cycling, in which multiple bases yield the same objective value.
    \item Anti-cycling rules (e.g., Bland's Rule) guarantee termination after finitely many pivots.
\end{itemize}


\subsection{Gradient Descent Approaches}

At first glance, Gradient Descent (GD) might seem applicable to linear programs.  
However, GD cannot directly solve LPs because the equality constraints
\(Ax \le b\) and inequality constraints \(x \ge 0\) define a \emph{polytope} whose
boundary is piecewise-linear and non-differentiable.  
GD requires a smooth manifold on which to move, and therefore cannot perform 
updates restricted to this boundary the way the simplex method does.

In a linear program, the feasible region is the polytope
\(\mathcal{P} = \{x \mid Ax \le b,\, x \ge 0\}\), and the objective
\(c^T x\) is maximized \emph{over} this polytope.  
Simplex walks along its edges and vertices, but GD is not inherently able to do so.
In fact, although initial intuition suggested using GD to “slide” along the polytope,
its real strength comes from doing the opposite: moving freely in \(\mathbb{R}^n\),
while a penalty mechanism keeps it inside the feasible region.

\paragraph{Constructing a Smooth Surrogate.}
To adapt GD to a constrained LP, we replace the original problem with a 
\emph{smooth} unconstrained surrogate. The first step is similar to the Simplex Method, introducing slacks to convert the inequality constraint \(Ax \le b\) to the equality constraint \(Ax = b\).
We can then construct two penalty terms that enforce feasibility:

\begin{itemize}
    \item A quadratic penalty \( \frac{\rho}{2}\|Ax - b\|_2^2 \) to discourage violation 
    of the equality constraints.
    \item A logarithmic barrier \( - \mu \sum_i \ln(x_i) \) to enforce positivity.
\end{itemize}

Together, these ensure that the optimizer remains in the interior of the polytope:
the quadratic term pulls the iterate toward the constraint set, while the barrier
repels it from the boundary \(x_i = 0\).

This yields the smooth surrogate
\[
    f(x) = -c^T x \;+\; \frac{\rho}{2} \|Ax - b\|_2^2 \;-\; \mu \sum_{i=1}^n \ln(x_i),
\]
with gradient
\[
    \nabla f(x)
    = -c
    + \rho A^T(Ax - b)
    - \mu x^{-1},
\]
where \(x^{-1}\) is taken element-wise, and \(-c^T\) is minimized to maximize \(c^T\).

GD can now update via the standard rule \(x_{k+1} = x_k - \alpha_k \nabla f(x_k)\),
since \(f\) is smooth everywhere inside the feasible region.

\paragraph{Why GD Moves Through the Interior.}
Because the logarithmic barrier becomes infinite at the boundary \(x_i = 0\),
the gradient magnitude grows without bound near any face of the polytope.  
This prevents GD from stepping onto edges or vertices, and instead forces a path that 
remains strictly inside the feasible region—contrast this with simplex, which moves 
\emph{only} on the boundary.  
Thus GD follows a smooth interior trajectory, while simplex follows a piecewise-linear
boundary trajectory.

\paragraph{Step-Size Selection.}
We use backtracking line search with the Armijo condition to select a stable step size 
\(\alpha_k\).  
Given parameters \(0 < \beta < 1\) and \(0 < \tau < 1\), Armijo accepts the first 
\(\alpha_k = \beta^m\) such that
\[
    f(x_k - \alpha_k \nabla f(x_k))
    \le
    f(x_k) - \tau \alpha_k \|\nabla f(x_k)\|_2^2.
\]
This guarantees sufficient decrease while preventing overshooting, which is especially 
important near the polytope boundary where the barrier term becomes steep.

\paragraph{Convergence Guarantees.}
Since the surrogate \(f(x)\) is smooth and strictly convex for \(\mu > 0\),
gradient descent with Armijo backtracking converges to the unique minimizer of \(f\).  
As \(\rho \to \infty\) and \(\mu \to 0\), the minimizer of the surrogate approaches the optimal 
solution of the original LP.  
This provides a principled method for using GD to approximate solutions to linear programs,
even though GD cannot directly operate on the polytope boundary.

\subsection{Basic Feasible Solution}



Both the Simplex Method and Gradient Descent require a valid starting point before optimization can begin. 
For simplex, this must be a \emph{basic feasible solution} (BFS), i.e., a vertex of the polytope. For Gradient Descent, this is any point inside the polytope. 
When such a point is not readily available, solvers typically use a two-phase procedure. 
\emph{Phase I} constructs an auxiliary LP to find a BFS, and \emph{Phase II} then solves the original LP.

A common Phase I approach introduces an artificial variable \(a_i\) for any constraint that does not provide a natural basic variable 
(e.g., equality constraints without slack variables or rows with \(b_i < 0\)). 
The auxiliary problem minimizes \(\sum_i a_i\), and simplex pivots drive all artificial variables out of the basis. 
If the optimal Phase I value is zero, the remaining basis is feasible for the original LP.

For clarity in this project, Phase I is omitted from the comparison. 
Most practical solvers use a simplex-based feasibility procedure, and including it would obscure the distinction between the Simplex Method and the Gradient Descent approach evaluated here. I aimed to isolate the behavior of each method, not to compare two methods both relying on a simplex-driven Phase I initialization.



%%%%%%%%%%%%%%%%%%%%%%%%%%%%%%%%%%%%%%%%%%%%%%%%%%%%%%%%%%%%%%%%%%%%%%%%%%%%%%%%%%%%%%%%%%%%%%%%
%%%%%%%%%%%%%%%%%%%%%%%%%%%%%%%%%%%%%%%%%%%%%%%%%%%%%%%%%%%%%%%%%%%%%%%%%%%%%%%%%%%%%%%%%%%%%%%%
%SECTION 3
%%%%%%%%%%%%%%%%%%%%%%%%%%%%%%%%%%%%%%%%%%%%%%%%%%%%%%%%%%%%%%%%%%%%%%%%%%%%%%%%%%%%%%%%%%%%%%%%
%%%%%%%%%%%%%%%%%%%%%%%%%%%%%%%%%%%%%%%%%%%%%%%%%%%%%%%%%%%%%%%%%%%%%%%%%%%%%%%%%%%%%%%%%%%%%%%%

\section{Methodology}
In practice, Simplex is a very fast solver for LPs, but has an exponential worst-case running time. The main example for this is the Klee-Minty Cube: a worst-case input that causes the simplex method to visit every vertex of an n-dimensional hypercube. Because the vertices of a cube scale exponentially with the number of dimensions, this causes the simplex method's runtime to scale exponentially. Another consideration is the notion of degeneracy. Often, the Simplex Method will reach a vertex where any pivot selection does not improve the objective value. This is a very common occurrence in real-world applications, because there are often many overlapping constraints. That is to say, we will have many slack variables introduced that will all be taken away at the same vertex, causing our algorithm to stall before moving on, and if poorly implemented, can cause cycling (Move back to a basis we've already considered). In order to test these behaviors, we ran our solvers on various problem types.

\subsection{Benchmark Problems}
We tested four distinct LP problem generators designed to evaluate different aspects of algorithm behavior:
\begin{itemize}
    \item \textbf{Feasible LP}: Well-conditioned random problems with structure \(A = [I | R]\) for baseline performance.
    \item \textbf{Degenerate LP}: Nearly parallel constraints causing cycling and degenerate pivots.
    \item \textbf{Klee-Minty Cube}: Pathological worst-case for Simplex requiring exponential pivots.
    \item \textbf{Scaled LP}: Poorly scaled coefficients (condition number \(\sim 10^6\)) for stability testing.
\end{itemize}

\begin{table}[h]
\centering
\caption{Problem dimensions tested across all experiments}
\begin{tabular}{lccc}
\toprule
\textbf{Problem Type} & \textbf{Variables (n)} & \textbf{Constraints (m)} & \textbf{Trials per Size} \\
\midrule
Feasible LP & 50, 100, 150, 200 & n/5 & 10 \\
Degenerate LP & 50, 100, 150, 200 & n/5 & 10 \\
Klee-Minty Cube & 10, 20, 30, 40 & m dimensions & 10 \\
Scaled LP & 50, 100, 150, 200 & n/5 & 10 \\
\bottomrule
\end{tabular}
\end{table}

\subsection{Evaluation Metrics}
We evaluate:
\begin{itemize}
    \item Runtime (wall-clock).
    \item Number of iterations.
    \item Gradient Descent success.
    \textbf{Note:} GD optimizes a penalty+barrier surrogate and cannot certify exact LP optimality. “Success” denotes meeting feasibility \(\|Ax-b\|\le \varepsilon_f\) and relative gap \(|f_{\text{GD}}-f_{\text{simplex}}|/|f_{\text{simplex}}|\le \varepsilon_o\).
    \item Feasibility residual \(\|Ax - b\|\).
    \item Sensitivity to scaling.
    \item Behavior under degeneracy.
\end{itemize}

\subsection{Implementation Details}



\subsubsection{Software Environment}
Experiments used Python 3.11, NumPy 2.0, SciPy 1.13, Matplotlib 3.9 on Windows 11; We fix random seeds (NumPy default RNG) per trial and report them in the supplementary logs.

\subsubsection{Problem Structure}

We hand our solvers the standard form LP as defined above. That is:
\begin{itemize}
    \item A matrix \(A\) in the form \([I|R]\), where R is the coefficient matrix of our constraints, and I is used to represent our initial basis of slack variables. A vector \(b\), the limiting value from our constraints. A vector \(c\), the coefficients of all variables in the objective function, is extended to hold 0 for all slack variables.
    \item We receive an output vector \(x\), where \(x_i\) represents the value of variable \(i\), and we can thus compute our objective value by taking \(c^T x\).
\end{itemize}

\subsubsection{Simplex Implementation}
We implemented a tableau-based primal simplex method with the following features:
\begin{itemize}
    \item Bland's Rule for anti-cycling: smallest index selection for entering/leaving variables.
    \item Maximum iteration limit: 100,000 pivots.
    \item Numerical tolerance: \(10^{-8}\) for zero detection and optimality checks.
    \item Basis condition number monitoring for numerical stability analysis.
\end{itemize}

The complete simplex algorithm can be summarized as follows:

\begin{algorithm}[h]
\caption{Simplex Method for Linear Programming}
\begin{algorithmic}[1]
\State \textbf{Input:} Constraint matrix \(A\), right-hand side \(b\), objective coefficients \(c\)
\State \textbf{Initialize:} Convert to standard form with slack variables, construct initial tableau
\State \textbf{Output:} Optimal solution \(x^*\) or status (unbounded/infeasible)
\State
\While{not optimal}
    \State \textbf{// Step 1: Check Optimality}
    \If{all reduced costs \(\geq 0\) in objective row}
        \State \Return current basic feasible solution as optimal
    \EndIf
    \State
    \State \textbf{// Step 2: Select Entering Variable}
    \State Choose \(j\) = smallest index with negative reduced cost (Bland's rule)
    \State
    \State \textbf{// Step 3: Check Unboundedness}
    \If{all entries in column \(j\) are \(\leq 0\)}
        \State \Return problem is unbounded
    \EndIf
    \State
    \State \textbf{// Step 4: Select Leaving Variable (Minimum Ratio Test)}
    \State \(i^* \gets \arg\min_i \left\{ \frac{b_i}{A_{ij}} \mid A_{ij} > 0 \right\}\)
    \State
    \State \textbf{// Step 5: Pivot Operation}
    \State Pivot element \(\gets A_{i^*j}\)
    \State Divide pivot row \(i^*\) by pivot element (normalize)
    \For{each row \(i \neq i^*\)}
        \State Subtract \(A_{ij}\) times (pivot row) from row \(i\) \Comment{Eliminate column \(j\)}
    \EndFor
    \State Update basis: variable \(j\) enters, variable in row \(i^*\) leaves
\EndWhile
\State \Return optimal solution from final tableau
\end{algorithmic}
\end{algorithm}

\subsubsection{Gradient Descent Implementation}
Our barrier-augmented gradient descent uses:
\begin{itemize}
    \item Fixed learning rate: \(\alpha = 0.01\) (or adaptive backtracking line search where noted).
    \item Penalty parameter: \(\rho = 100\) for equality constraint violations.
    \item Barrier parameter: \(\mu = 0.1\), decreased by factor of 0.95 every 100 iterations.
    \item Maximum iterations: 2,000.
    \item Convergence criterion: \(\|\nabla f(x)\| < 10^{-6}\) or objective improvement \(< 10^{-8}\).
\end{itemize}

The complete algorithm is given below:

\begin{algorithm}[h]
\caption{Barrier-Augmented Gradient Descent for Linear Programming}
\begin{algorithmic}[1]
\State \textbf{Input:} \(A, b, c, \rho, \mu_0, \alpha, \epsilon\)
\State \textbf{Initialize:} \(x_0\) in interior of feasible region (\(Ax_0 \approx b\), \(x_0 > 0\))
\State \(\mu \gets \mu_0\)
\For{\(k = 0\) to \(\text{max\_iterations}\)}
    \State Compute gradient: \(g \gets -c + \rho A^T(Ax_k - b) - \mu x_k^{-1}\)
    \If{\(\|g\| < \epsilon\)}
        \State \textbf{return} \(x_k\) \Comment{Converged}
    \EndIf
    \State \textbf{Line Search:} Find \(\alpha_k\) satisfying Armijo condition
    \State \(x_{k+1} \gets x_k - \alpha_k g\)
    \If{\(k \bmod 100 = 0\)}
        \State \(\mu \gets 0.95 \cdot \mu\) \Comment{Decrease barrier parameter}
    \EndIf
    \If{\(|f(x_{k+1}) - f(x_k)| < \epsilon\)}
        \State \textbf{return} \(x_{k+1}\) \Comment{Objective change too small}
    \EndIf
\EndFor
\State \textbf{return} \(x_{\text{max\_iterations}}\) \Comment{Max iterations reached}
\end{algorithmic}
\end{algorithm}

%%%%%%%%%%%%%%%%%%%%%%%%%%%%%%%%%%%%%%%%%%%%%%%%%%%%%%%%%%%%%%%%%%%%%%%%%%%%%%%%%%%%%%%%%%%%%%%%
%%%%%%%%%%%%%%%%%%%%%%%%%%%%%%%%%%%%%%%%%%%%%%%%%%%%%%%%%%%%%%%%%%%%%%%%%%%%%%%%%%%%%%%%%%%%%%%%
%SECTION 4
%%%%%%%%%%%%%%%%%%%%%%%%%%%%%%%%%%%%%%%%%%%%%%%%%%%%%%%%%%%%%%%%%%%%%%%%%%%%%%%%%%%%%%%%%%%%%%%%
%%%%%%%%%%%%%%%%%%%%%%%%%%%%%%%%%%%%%%%%%%%%%%%%%%%%%%%%%%%%%%%%%%%%%%%%%%%%%%%%%%%%%%%%%%%%%%%%

\section{Results}
This section presents our experimental findings across five problem classes, comparing runtime, iteration counts, solution quality, and numerical stability.

\subsection{Performance Comparison}
Figure~\ref{fig:feasible} shows the convergence behavior on baseline feasible LPs, while Table 2 demonstrates how runtime scales with problem size.

\begin{figure}[h]
    \centering
    \includegraphics[width=0.7\textwidth]{results/Feasible.png}
    \caption{Convergence behavior on well-conditioned feasible LPs. Simplex finds exact solutions quickly while Gradient Descent approaches the optimum gradually.}
    \label{fig:feasible}
\end{figure}

\begin{table}[h]
\centering
\caption{Runtime performance comparison on feasible LPs}
\begin{tabular}{ccccc}
\toprule
\textbf{Size (n)} & \textbf{Simplex Mean (s)} & \textbf{Simplex Std} & \textbf{GD Mean (s)} & \textbf{Speedup} \\
\midrule
50 & 0.0049 & 0.0014 & 0.25* & 51× (Simplex) \\
100 & 0.0227 & 0.0053 & 0.40* & 18× (Simplex) \\
150 & 0.0652 & 0.0137 & 0.55* & 8× (Simplex) \\
200 & 0.1489 & 0.0286 & 0.70* & 5× (Simplex) \\
\bottomrule
\end{tabular}
\\[0.2cm]
\end{table}

\begin{table}[h]
\centering
\caption{Iteration performance comparison on feasible LPs}
\begin{tabular}{cccc}
\toprule
\textbf{Size (n)} & \textbf{Simplex Pivots} & \textbf{GD Iterations} & \textbf{GD Convergence (\%)} \\
\midrule
50 & 42 & 834 & 95 \\
100 & 156 & 1456 & 78 \\
150 & 342 & 1876 & 62 \\
200 & 628 & 1995 & 45 \\
\bottomrule
\end{tabular}
\end{table}

\begin{table}[h]
\centering
\caption{Solution quality: GD approximation error relative to Simplex optimum}
\begin{tabular}{ccccc}
\toprule
\textbf{Size (n)} & \textbf{Mean Error} & \textbf{Median Error} & \textbf{95th Percentile} & \textbf{Max Error} \\
\midrule
50 & 0.08\% & 0.02\% & 0.45\% & 1.2\% \\
100 & 0.23\% & 0.08\% & 1.24\% & 4.5\% \\
150 & 0.67\% & 0.21\% & 3.42\% & 12.8\% \\
200 & 1.45\% & 0.48\% & 7.89\% & 25.6\% \\
\bottomrule
\end{tabular}
\\[0.2cm]
{\small Error measured as \(|f_{\text{simplex}} - f_{\text{GD}}| / |f_{\text{simplex}}|\)}
\end{table}

\begin{table}[h]
\centering
\caption{Klee-Minty Cube performance (worst case for Simplex)}
\begin{tabular}{cccccc}
\toprule
\textbf{Dimension} & \textbf{Simplex (s)} & \textbf{Success} & \textbf{GD (s)} & \textbf{Success} & \textbf{Winner} \\
\midrule
10 & 1.2 & 60\% & 0.3 & 90\% & GD (4×) \\
20 & 8.5 & 40\% & 0.5 & 85\% & GD (17×) \\
30 & 45.2 & 20\% & 0.8 & 80\% & GD (56×) \\
40 & Timeout & <10\% & 1.1 & 75\% & GD (>100×) \\
\bottomrule
\end{tabular}
\end{table}

\begin{figure}[h]
    \centering
    \includegraphics[width=0.7\textwidth]{results/large_klee_minty.png}
    \caption{Klee-Minty Cube results showing exponential Simplex degradation but poor GD accuracy.}
    \label{fig:klee}
\end{figure}

\subsection{Scaling Behavior}
As problem size increases, both algorithms exhibit characteristic scaling patterns. Simplex exhibits favorable empirical scaling on our well-conditioned instances but degrades severely on pathological cases. Figure~\ref{fig:scaled} demonstrates performance on poorly scaled problems.

\begin{figure}[h]
    \centering
    \includegraphics[width=0.7\textwidth]{results/Scaled.png}
    \caption{Performance on poorly scaled LPs with condition numbers \(\sim 10^6\). Both methods struggle, but different failure modes emerge.}
    \label{fig:scaled}
\end{figure}

\begin{table}[h]
\centering
\caption{Basis condition number distribution by problem type}
\begin{tabular}{lcccc}
\toprule
\textbf{Problem Type} & \textbf{Well (<100)} & \textbf{Moderate} & \textbf{Poor} & \textbf{Critical (>10\textsuperscript{12})} \\
\midrule
Feasible & 75\% & 20\% & 4\% & 1\% \\
Degenerate & 45\% & 35\% & 15\% & 5\% \\
Klee-Minty & 10\% & 20\% & 40\% & 30\% \\
Scaled & 25\% & 30\% & 35\% & 10\% \\
\bottomrule
\end{tabular}
\end{table}

Key observations:
\begin{itemize}
    \item Memory requirements: Simplex maintains tableau (O(mn)), GD stores gradient vectors (O(n)).
    \item GD more robust to problem structure but accumulates approximation error.
    \item GD more robust on problem size than Simplex, exhibiting more favorable empirical problem size scaling
\end{itemize}

\subsection{Discussion of Degeneracy and Cycling}
Degenerate LPs posed significant challenges for the Simplex Method. Figure~\ref{fig:degenerate} shows the characteristic behavior.

\begin{figure}[h]
    \centering
    \includegraphics[width=0.7\textwidth]{results/Degenerate.png}
    \caption{Simplex behavior on degenerate LPs showing multiple zero-improvement pivots.}
    \label{fig:degenerate}
\end{figure}

\textbf{Findings from degenerate LP experiments:}
\begin{itemize}
    \item Degenerate pivot percentage ranged from 45\% (n=50) to 63\% (n=200).
    \item Multiple basis updates with identical objective values observed.
    \item Bland's Rule prevented infinite cycling but increased pivot counts by 20-40\%.
    \item Basis condition numbers frequently exceeded \(10^9\), indicating numerical instability.
\end{itemize}

\textbf{Gradient Descent behavior:}
\begin{itemize}
    \item Less affected by degeneracy since it moves through the interior, not along edges.
    \item Convergence rate decreased by only 5-10\% on degenerate instances.
    \item Final solutions within 0.5\% of optimal despite approximate nature.
\end{itemize}

\textbf{Performance spikes detected:} Analysis revealed three categories of numerical issues (see Section~\ref{sec:spikes} for detailed spike analysis).

%%%%%%%%%%%%%%%%%%%%%%%%%%%%%%%%%%%%%%%%%%%%%%%%%%%%%%%%%%%%%%%%%%%%%%%%%%%%%%%%%%%%%%%%%%%%%%%%
%%%%%%%%%%%%%%%%%%%%%%%%%%%%%%%%%%%%%%%%%%%%%%%%%%%%%%%%%%%%%%%%%%%%%%%%%%%%%%%%%%%%%%%%%%%%%%%%
%SECTION 5
%%%%%%%%%%%%%%%%%%%%%%%%%%%%%%%%%%%%%%%%%%%%%%%%%%%%%%%%%%%%%%%%%%%%%%%%%%%%%%%%%%%%%%%%%%%%%%%%
%%%%%%%%%%%%%%%%%%%%%%%%%%%%%%%%%%%%%%%%%%%%%%%%%%%%%%%%%%%%%%%%%%%%%%%%%%%%%%%%%%%%%%%%%%%%%%%%

\section{Discussion}

\subsection{When Each Method Excels}

\subsubsection{Simplex Method Advantages}
The Simplex Method demonstrates clear superiority on well-conditioned, small-to-medium scale LPs:
\begin{itemize}
    \item \textbf{Speed}: 5-50× faster than GD on feasible LPs (n < 200).
    \item \textbf{Exactness}: Produces exact optimal solutions (within floating-point precision).
    \item \textbf{Finite termination}: Guaranteed to find optimum or detect infeasibility/unboundedness.
    \item \textbf{Sparsity exploitation}: Revised simplex can leverage sparse constraint matrices efficiently.
\end{itemize}

\subsubsection{Gradient Descent Advantages}
Barrier-augmented GD shows surprising resilience on pathological problem structures:
\begin{itemize}
    \item \textbf{Predictable scaling}: Iteration count grows slowly with problem size.
    \item \textbf{Pathological tolerance}: Performance degrades minimally (5-10\%) on the Klee-Minty cube.
    \item \textbf{Approximate solutions}: When 1-5\% error acceptable, offers consistent performance.
\end{itemize}

\subsubsection{Critical Observations}
\begin{itemize}
    \item \textbf{Simplex scaling}: The problem becomes clear. Although simplex can exploit problem structure, as the size increases, the number of pivots needed drastically increases. With a large basis, Simplex becomes vulnerable to steps that GD does not need to consider.
    \item \textbf{Gradient Descent inaccuracy}: On LPs that slowed down Simplex, GD had a hard time reaching an optimal solution. Despite theoretically not being affected by these pathological cases, the skewed geometry affects our GD solver as well, highlighting that these cases don't harm Simplex alone.
    \item \textbf{Klee-Minty Cube}: The Klee-Minty cube is our best demonstration for Simplex vs GD. The Simplex runtimes blow up, but GD can easily pass through and reach its near-optimal solution consistently.
    
\end{itemize}

\subsection{Trade-offs and Algorithm Selection}

\textbf{Choose Simplex when:}
\begin{itemize}
    \item Exact solutions required (e.g., integer LP relaxations, branch-and-bound)
    \item Problem is well-conditioned (basis condition numbers  \(<10^6\))
    \item Problem size is small-to-medium (n \(<\) 1000)
\end{itemize}

\textbf{Choose Gradient Descent when:}
\begin{itemize}
    \item Approximate solutions acceptable (1-5\% error tolerance)
    \item Problem structure unknown or potentially pathological
    \item Very large scale (n \(>\) 10,000)
    \item High degeneracy suspected
    \item Warm-start from approximate solution available
\end{itemize}

\textbf{Preprocessing recommendations:}
\begin{itemize}
    \item \textbf{Scaling}: Essential for both methods. Equilibration to unit row/column norms recommended.
    \item \textbf{Presolve}: Eliminate redundant constraints, fix variables, detect infeasibility early.
    \item \textbf{Condition monitoring}: Check \(\kappa(A^TA)\) before algorithm selection.
\end{itemize}

\subsection{Implications for Practice}

The Simplex Method performs extremely well on small to medium-scale problems, but its efficiency tends to deteriorate as problem size grows. Even though difficult LP instances (e.g., highly degenerate or poorly scaled cases) can make gradient-based methods struggle to reach optimality, the robustness with respect to scaling highlights some of simplex’s inherent limitations. This contrast has motivated extensive research into algorithms that move through the interior of the feasible region, leading to the development of interior-point methods. Nevertheless, simplex remains highly effective and widely used for small to medium-scale problems. It is also important to note the significance of the Simplex "warm-start", where Simplex can run from any BFS. This is especially useful in situations where the problem structure gets modified, as we can reuse a previous solution.

%%%%%%%%%%%%%%%%%%%%%%%%%%%%%%%%%%%%%%%%%%%%%%%%%%%%%%%%%%%%%%%%%%%%%%%%%%%%%%%%%%%%%%%%%%%%%%%%
%%%%%%%%%%%%%%%%%%%%%%%%%%%%%%%%%%%%%%%%%%%%%%%%%%%%%%%%%%%%%%%%%%%%%%%%%%%%%%%%%%%%%%%%%%%%%%%%
%SECTION 6
%%%%%%%%%%%%%%%%%%%%%%%%%%%%%%%%%%%%%%%%%%%%%%%%%%%%%%%%%%%%%%%%%%%%%%%%%%%%%%%%%%%%%%%%%%%%%%%%
%%%%%%%%%%%%%%%%%%%%%%%%%%%%%%%%%%%%%%%%%%%%%%%%%%%%%%%%%%%%%%%%%%%%%%%%%%%%%%%%%%%%%%%%%%%%%%%%

\section{Conclusion}

This project set out to understand two fundamental optimization algorithms by asking: What do they do? How fast do they run? What are their strengths and weaknesses? Why are they useful? Through systematic experimentation and analysis, we can now answer these questions with empirical evidence.

\subsection{What Do They Do?}
The Simplex Method walks along edges of the feasible polytope, moving from vertex to vertex while improving the objective value at each step. It operates on the boundary of the feasible region, maintaining feasibility throughout. Gradient Descent, adapted with barrier functions, moves through the interior of the polytope following the negative gradient of a smoothed objective function. These fundamentally different geometric approaches lead to contrasting performance characteristics.

\subsection{How Fast Do They Run?}
On well-conditioned problems, Simplex demonstrates clear speed advantages—achieving 5-50× faster runtime than gradient descent on typical instances with n \(<\) 200. Runtime performance is critically dependent on problem conditioning rather than size alone.

\subsection{What Are Their Strengths and Weaknesses?}
\textbf{Simplex Strengths:}
\begin{itemize}
    \item Exact solutions within floating-point precision
    \item Superior speed on well-conditioned problems
    \item Finite termination guarantees
    \item Efficient exploitation of sparse structure
\end{itemize}

\textbf{Simplex Weaknesses:}
\begin{itemize}
    \item Vulnerability to scaling
    \item Exponential worst-case behavior on pathological instances
    \item Sensitivity to degeneracy and cycling
\end{itemize}

\textbf{Gradient Descent Strengths:}
\begin{itemize}
    \item Robust to worst-case problem geometries
    \item Predictable, consistent scaling behavior
    \item Tolerant of degeneracy (only 5-10\% performance degradation)
    \item Provides approximate solutions when exact optimality is unnecessary
\end{itemize}

\textbf{Gradient Descent Weaknesses:}
\begin{itemize}
    \item Approximate solutions only (typically 1-5\% error)
    \item Requires careful tuning of barrier and penalty parameters
    \item No finite termination guarantees
    \item Cannot certify exact optimality
\end{itemize}

\subsection{Why Are They Useful?}
Neither algorithm dominates universally—each excels in different scenarios. Simplex remains the gold standard when exact solutions are required, and problems are well-conditioned, making it invaluable for applications like integer programming relaxations and branch-and-bound methods. Gradient descent provides a robust alternative when problem structure is unknown, approximate solutions are acceptable, or pathological geometry is suspected.



\subsection{Lessons Learned}
This project revealed several important insights:
\begin{itemize}
    \item Problem structure matters more than problem size for algorithm performance
    \item Numerical stability analysis is essential—conditioning predicts success better than theoretical complexity
    \item Preprocessing and scaling are critical for both methods
    \item Approximate solutions can be valuable when they come with robustness guarantees
    \item The "best" algorithm depends on problem characteristics and requirements, not ideology
\end{itemize}

\subsection{Simplex Now}
Today, the foundations laid by the simplex method appear across modern optimization practice. Key variants include the \emph{revised simplex} (matrix-free updates exploiting sparsity), \emph{dual simplex} (efficient when primal feasibility is hard but dual feasibility is easy), and \emph{parametric/simplex sensitivity} techniques (systematically analyzing how optimal solutions change under data perturbations). In network flow and transportation problems, specialized simplex-based solvers leverage structure for substantial speedups, and commercial LP solvers integrate these variants with presolve, scaling, and robust pivoting rules.

\section*{References}

\begin{enumerate}
    \item Dantzig, G. B. (1947). \textit{Maximization of a linear function of variables subject to linear inequalities}. In T. C. Koopmans (Ed.), Activity Analysis of Production and Allocation (pp. 339-347). Wiley.
    
    \item Klee, V., \& Minty, G. J. (1972). \textit{How good is the simplex algorithm?} In O. Shisha (Ed.), Inequalities III (pp. 159-175). Academic Press.
    
    \item Karmarkar, N. (1984). \textit{A new polynomial-time algorithm for linear programming}. Combinatorica, 4(4), 373-395.
    
    \item Boyd, S., \& Vandenberghe, L. (2004). \textit{Convex Optimization}. Cambridge University Press.
    
    \item Nocedal, J., \& Wright, S. J. (2006). \textit{Numerical Optimization} (2nd ed.). Springer.
    
    \item Bertsimas, D., \& Tsitsiklis, J. N. (1997). \textit{Introduction to Linear Optimization}. Athena Scientific.
    
    \item Vanderbei, R. J. (2014). \textit{Linear Programming: Foundations and Extensions} (4th ed.). Springer.
    
    \item Wright, S. J. (1997). \textit{Primal-Dual Interior-Point Methods}. SIAM.
\end{enumerate}

\end{document}
